\section{Actual Type I points: host automorphy and adjoint rigidity}
\label{sec:cubic-typei-host-automorphy-rigidity}

The support conclusion for the actual tuple has a characteristic-zero
application as well as the curve application of AC-G109. It gives
automorphy over the chosen CM host. Adjoint Selmer vanishing then
descends by restriction and corestriction. This last descent is a
cohomological statement; it does not identify an automorphic
representation over the original field with the level required for
the counting theorem.

\paragraph{Actual data and retained support.}
Put \(p=269\), \(K=\mathbb Q(\omega)\), and use the actual fixed-type
model \(\mathscr S\) of AC-G65--AC-G66 over its original coefficient
DVR \(\mathcal O_0\). All its original datum and source hypotheses
are retained. Write \(c_p\) for the actual cyclotomic character,
so the fixed ordered inertia characters are
\[
       \chi^j\psi c_p^e,\quad 1,\quad
                   \chi^j\psi^{-1}c_p^{-e},\qquad e\ge1.
 \tag{G110.1}
\]
Their residual tuple is denoted by \(\boldsymbol\theta\).
In particular the retained datum has \(268\nmid2e\), and
the local separating element of AC-G65-flag exists.

Fix one actual host \(H/F\), \(F=H^+\), of AC-G93--AC-G108,
and its coefficient DVR \(\mathcal O\), residue field \(k\),
and rings \(B\subset A\subset R\). The native ordinary
determinant ring is written \(\mathcal B\), with native map
\(\Phi:\mathcal B\to B\), as in AC-G97. Use the actual
native support ideal \(J\subset\mathcal B\) and retain its
\texttt{Modular} properties, including arithmetic-point
classicality. The precise support condition for this section is
\[
 \Phi(J)\subset\mathfrak q\cap B
       \quad\hbox{for every actual Type I prime }
                    \mathfrak q\subset A.
 \tag{G110.2}
\]
This is the Type I conclusion of AC-G108 for the image of the
native ideal. All premises used to obtain that conclusion,
including uniform endpoint support \((E)\), remain premises
until supplied. A finite-weight ideal unrelated to \(J\) is
not a replacement for (G110.2).

The two coefficient DVRs may differ. They and the original
algebraic point are embedded in the fixed algebraic closure
\(\overline{\mathbb Q}_p\); their residual embeddings identify
the already specified tuple. No unramifiedness of the large
native coefficient ring is inferred, and the original
\(e_{\mathcal O_0}=1\) condition, where imposed, is unchanged.

\begin{lemma}[AC-G110-irr: irreducibility on every open subgroup and on the cyclotomic kernel]
\label{AC-G110-irr}
Let \(f:\mathscr S\to\overline{\mathbb Q}_p\) be an actual
algebraic local point of Type I, and let
\(\rho:G_K\to\operatorname{GL}_3(\overline{\mathbb Q}_p)\)
be its multiplier-one restriction. Then \(\rho\) is absolutely
irreducible on every open subgroup of \(G_K\). For every finite
extension \(H'/K\), its restriction to
\[
                G_{H'(\mu_{p^\infty})}
                       =G_{H'}\cap\ker c_p
 \tag{G110.3}
\]
is absolutely irreducible as well.
\end{lemma}
\begin{proof}
The finite-image and continuity argument of AC-G58 and AC-G66
places all matrices in the integers of one finite extension of
\(\mathbb Q_p\), with exactly the residual diagonal tuple of
the original model. The representation is irreducible over
the algebraic closure by the definition of Type I.

First let \(U\) be an open normal subgroup of \(G_K\).
Apply the generic prime-rank theorem \texttt{irrNormal} in
\path{CurveDVRSpecialization/IrrNormal.lean}. Its integral
reduction datum is the just specified matrix representation
and residue map. Take its local group to be the full group
at \(p\), with the ordered flag characters of AC-G65-flag.
The required factorization of local characteristic polynomials
follows from that flag. A nonscalar residual local witness is
the element in (RC.17), for which the three residual diagonal
values are distinct.

For its remaining hypothesis, the intersection of \(U\) with
local inertia is open in inertia. Intersect further with the
kernels of the finite characters \(\chi,\psi\). On this open
subgroup every ratio of two different characters in (G110.1)
is a nonzero power \(c_p^{\pm e}\) or \(c_p^{\pm2e}\).
The cyclotomic image of an open inertia subgroup is infinite;
none of these ratios is identically one. These are precisely
all hypotheses of \texttt{irrNormal}, with prime rank three,
so \(\rho|_U\) is irreducible. Taking the normal core of any
open subgroup proves the first assertion.

Here is the separate argument for the infinite cyclotomic
extension; the preceding local-ratio argument is not applied
to its kernel. Work on \(G'=G_{H'}\) and put
\(N=G'\cap\ker c_p\). Every open subgroup of \(G'\) is
open in \(G_K\), so the result just proved still holds.
The prime-rank Clifford alternatives for a normal subgroup
say that \(\rho|_N\) is irreducible, is scalar, or is a sum
of three distinct character lines permuted transitively by
\(G'\).

For the algebraic content of this assertion, choose a simple
nonzero \(N\)-submodule of the three-dimensional space.
The sum of its \(G'\)-translates is the whole space by
irreducibility; it is a sum of simple \(N\)-modules, hence
semisimple. Its isotypic components are permuted transitively
and have equal dimensions. Their number divides three.
Three components give three distinct lines. One component
is a sum of copies of one simple module, whose dimension
divides three; it is either the entire irreducible module or
one character acting by scalars. This proves the alternatives
also for the closed, generally nonopen subgroup \(N\).

In the three-line case the kernel of the permutation action
is closed and has finite index: the stabilizer of each line
is given by finitely many continuous linear equations in the
matrix entries. All needed basis coefficients lie in one
finite extension. A closed subgroup of finite index in a
profinite group is open. The representation on that kernel
preserves a line, contradicting the first assertion.

In the scalar case \(G'/N\) is the compact cyclotomic image,
a closed subgroup of \(\mathbb Z_p^\times\). Since \(p\)
is odd, \(\mathbb Z_p^\times\) is procyclic, as is each
closed subgroup: its finite prime-to-\(p\) cyclic factor
and its pro-\(p\) cyclic factor have relatively prime
orders at every finite stage. Choose \(t\in G'\) lifting
a topological generator of that image, and an eigenline
of \(\rho(t)\). Its closed stabilizer contains \(N\)
and \(t\). Its image in \(G'/N\) is compact, hence
closed, and contains a dense cyclic subgroup. It is the
whole quotient, so the stabilizer is all of \(G'\).
This contradicts irreducibility of rank three. Only the
irreducible alternative remains, proving (G110.3).
\end{proof}

\begin{theorem}[AC-G110-point: every actual Type I point is automorphic on the support host]
\label{AC-G110-point}
Under the actual setup and (G110.2), for every such point \(f\)
the representation
\[
               \nu=\rho|_{G_H}\otimes c_p^{-1}
 \tag{G110.4}
\]
satisfies the literal native \texttt{Af.IsAutCsd} condition:
it is attached to a regular algebraic, cuspidal, conjugate
self-dual automorphic representation over \(H\).
Its restriction to \(G_{H(\mu_{p^\infty})}\) is irreducible.
\end{theorem}
\begin{proof}
\emph{1. A test ring with the native residue field.}
Choose a finite extension \(E'/\mathbb Q_p\) containing the
matrix coefficients of the point and the native coefficient
field. All matrices and their inverses are integral by the
finite-image result already used. Write \(D=\mathcal O_{E'}\)
and let \(k'\) be its residue field. The prescribed embeddings
identify \(k\) with a subfield of \(k'\). Put
\[
               C=\{x\in D:x\bmod\mathfrak m_D\in k\}.
 \tag{G110.5}
\]
It contains \(\mathcal O\), is an \(\mathcal O\)-submodule
of the finite module \(D\), and is therefore finite and
Noetherian. It is local, its maximal ideal is
\(\mathfrak m_D\), its residue field is exactly \(k\), and
\(\mathfrak m_C^n=\mathfrak m_D^n\). The inverse image of
a finite subset of the residue field is closed in \(D\),
so it is complete for this topology. The restricted
representation and its inverse matrices have their residual
entries in \(k\), hence have coefficients in \(C\).
This transfers the representation only; no map of the
original deformation ring over a changed coefficient DVR
is asserted.

\emph{2. The actual host point.}
The original killing condition annihilates all away-inertia
generators in the actual host joint kernel. Its continuous
kernel therefore contains their closed normal closure, so
the representation descends to the group of AC-G94.
Its exact polarization has form matrix \(J\).
At the host witness form
\(Q=\rho(g_0)-\rho(g_0^\star)\).
It is \(J\)-skew, with characteristic polynomial
\(X(X^2-b)\), where \(b=\operatorname{Tr}(Q^2)/2\)
reduces to the square of the nonzero residual difference.
The same Hensel projectors and column rescalings as in
AC-G109-curve give a strict \(J\)-orthogonal matrix over
\(C\) diagonalizing \(Q\) to \((a,0,-a)\).
The fixed local flags restrict to every full host local
group. Their labeled integral projectors define them and
their quotient characters over \(C\); strict normalization
preserves the ordered-product identities and those characters.
These are all conditions of the actual AC-G94 functor.
It supplies a map \(R\to C\), and hence
\[
              \xi:A\longrightarrow C
                       \hookrightarrow\overline{\mathbb Q}_p.
 \tag{G110.6}
\]
AC-G110-irr makes the restricted representation irreducible,
so \(\ker\xi\) is an actual Type I prime. Equation
(G110.2) therefore makes
\[
              x=\xi|_B\circ\Phi:
                 \mathcal B\longrightarrow\overline{\mathbb Q}_p
                    \quad\hbox{kill }J.
 \tag{G110.7}
\]

\emph{3. Check the exact arithmetic-point input.}
The native Pair/map identities of AC-G97 identify the
determinant at \(x\), after pullback from the native group,
with \(\det\nu\). This is equality of determinant laws,
not only equality of traces at group elements. They also
identify all full local characters with the original
ordered characters multiplied by \(c_p^{-1}\).
Thus \(x.\texttt{IsRep}\,\nu\) holds.

On coefficient-prime inertia these local characters are
\[
       \chi^j\psi c_p^{e-1},\quad c_p^{-1},\quad
                     \chi^j\psi^{-1}c_p^{-e-1}.
 \tag{G110.8}
\]
In the native geometric local Artin convention, a finite
character times \(c_p^m\) has locally algebraic weight
\(-m\) at every embedding. This is the literal
\texttt{laCyc}/finite-character calculation contained in
\texttt{AutPtOf.ap\_locAlg}; the finite factors in (G110.8)
have open kernel and are not discarded as characters.
Consequently the native weights are
\[
                         (1-e,\ 1,\ 1+e),
 \tag{G110.9}
\]
strictly increasing because \(e\ge1\). Hence
\(x.\texttt{Arith}\) holds on all required local groups.
The representation \(\nu\) is irreducible by AC-G110-irr.

Apply exactly \texttt{arithmeticPointClassicality}, the
\texttt{end2} field of the retained \texttt{HidaProps} in
\path{WeilGroupCharacters/ArithmeticPointClassicality.lean}.
Its premises are \(J\subset\ker x\), \texttt{Arith},
\texttt{IsRep}, and irreducibility, all now checked.
Its conclusion is the asserted \texttt{Af.IsAutCsd} for
\(\nu\). This use is not the \texttt{PT} direction,
which starts with an already automorphic point.
Finally scalar twisting preserves invariant subspaces,
so (G110.3) gives the asserted cyclotomic irreducibility.
\end{proof}

\begin{theorem}[AC-G110-rigidity: original rational adjoint vanishing at Type I points]
\label{AC-G110-rigidity}
Under the same hypotheses, write \(r_f\) for the original
semidirect representation of \(G_{\mathbb Q}\) at \(f\),
and \(\operatorname{ad}(r_f)\) for its actual linear
adjoint on \(M_3(\overline{\mathbb Q}_p)\). Then
\[
        H^1_f\bigl(\mathbb Q,\operatorname{ad}(r_f)\bigr)=0.
 \tag{G110.10}
\]
In particular every actual dual-point tangent at this Type I
point is a coboundary. This supplies the Type I part of the
original point-rigidity condition at \(269\).
\end{theorem}
\begin{proof}
We use the published theorem of J.~A.~Thorne,
\emph{On the vanishing of adjoint Bloch--Kato Selmer groups
of irreducible automorphic Galois representations},
\emph{Pure Appl. Math. Q.}~18 (2022), no.~5, 2159--2202,
Theorem~A, precisely Theorem~6.1. It gives vanishing over
the maximal real subfield for the polarized adjoint of a
regular algebraic cuspidal polarizable representation
whose Galois representation stays irreducible on the full
cyclotomic kernel. The primary author version is
\url{https://www.dpmms.cam.ac.uk/~jat58/adequate_new.pdf},
Theorem~6.1 on printed p.~23; arXiv:2207.04925.
This published theorem is retained as an input; no new
proof of its patching argument is claimed.

AC-G110-point supplies precisely such an automorphic
representation on \(H\), with associated representation
\(\nu\). Native \texttt{AttC} specifies the same normalized
unramified characteristic polynomials, and semisimplicity;
Chebotarev and characteristic-zero Brauer--Nesbitt identify
it with the associated Galois representation. The exact
cyclotomic irreducibility is (G110.3), not a residual
irreducibility assertion. Thorne's theorem gives
\[
                    H^1_f(F,\operatorname{Ad}\nu)=0.
 \tag{G110.11}
\]

We check the adjoint, including its action on the other
CM coset. On \(G_H\), scalar twisting by \(c_p^{-1}\)
does not change conjugation on matrices. The original
polarization is
\(\rho(cgc^{-1})=J\rho(g)^{-t}J\), so
\[
             \nu(cgc^{-1})
                    =c_p(g)^{-2}J\nu(g)^{-t}J.
 \tag{G110.12}
\]
It is the rank-three cyclotomic normalization, with
symmetric form \(J\). The polarized adjoint over \(F\)
therefore has the action \(X\mapsto-JX^tJ\) at the
chosen complex conjugation. This is exactly the
differential of the original semidirect inverse-transpose
involution, because \(r_f(c).\mathrm{mat}=I\).
If the automorphic representation is identified in a
different basis, irreducibility makes its polarization
unique up to scalar, which leaves the adjoint unchanged.
Thus \(\operatorname{Ad}\nu\) is the restriction of
\(\operatorname{ad}(r_f)\) to \(G_F\), with the actual
involution, not the unrelated untwisted matrix module.

Restriction carries Bloch--Kato classes from \(\mathbb Q\)
to \(F\): away-inertia coboundaries remain such after
restriction, and the crystalline-period local kernel
restricts to the corresponding kernel at every place
above \(p\). On rational continuous cohomology,
restriction followed by corestriction is multiplication
by \([F:\mathbb Q]\), an invertible scalar in the
coefficient field. Hence restriction is injective,
including when the degree is divisible by \(p\).
Equation (G110.11) gives (G110.10).

Finally AC-G66-BK places every actual algebraic dual-point
tangent in this exact Bloch--Kato group. Its class is
zero by (G110.10), which is the asserted coboundary
condition. No claim that every coordinate derivative
itself vanishes is needed or made.
\end{proof}

\paragraph{Scope for the main arithmetic bound.}
These conclusions use the actual Type I support (G110.2)
and the retained native classicality input. The uniform
endpoint ideal needed in AC-G108 is a separate supply
obligation until its proof is completed. Type II at
\(269\) still has the mixed coefficient and actual
tangent-image obstruction of AC-G68; this theorem does
not remove it. Accordingly full \texttt{M4}, the model
length estimate, and H04.4 are not conclusions here.
The automorphy conclusion is over \(H\). It does not
by itself supply the original-\(K\) automorphic
descent, \texttt{WitA}, \texttt{Level}, or \texttt{Unl}
data of the fixed-field integer-product counting theorem.

\paragraph{Exact source and geometric dictionary.}
The native files used here are at the same pinned commit
\texttt{7cbf7a0c32fca1badf05931879a37c3f8a2b03b4}.
Complete containing reads include \path{IrrNormal.lean},
\path{LinClosedOfAdic.lean}, \path{StrongIrrOfOpen.lean},
\path{AutPtOf.lean},
\path{WeilGroupCharacters/ArithmeticPointClassicality.lean},
and \path{AutomorphicForms/BlochKatoSelmerVanishes.lean}.
Their exact paths, hashes and scopes are in the source
receipt. No old one-character \texttt{RigidityFirstCase}
or \texttt{AutEnd3} wrapper is substituted for the
actual cubic point. No Lean execution is asserted.

Geometrically the integral representation selects a point
of the actual ordinary deformation chart; the invariant
support ideal places its determinant on the automorphic
support. The final result is the vanishing of a rational
adjoint deformation space, descending along a finite
cover. It is not an integral bound on its torsion
or a construction of an absolute intersection pairing.
The finite-cover language uses the already audited
SGA~1 dictionary without a new archival attribution.

